\documentclass[10pt, letterpaper]{article}
 
% Packages:
\usepackage[
    ignoreheadfoot, % set margins without considering header and footer
    top=2 cm, % seperation between body and page edge from the top
    bottom=2 cm, % seperation between body and page edge from the bottom
    left=2 cm, % seperation between body and page edge from the left
    right=2 cm, % seperation between body and page edge from the right
    footskip=1.0 cm, % seperation between body and footer
    % showframe % for debugging 
]{geometry} % for adjusting page geometry
\usepackage{titlesec} % for customizing section titles
\usepackage{tabularx} % for making tables with fixed width columns
\usepackage{array} % tabularx requires this
\usepackage[dvipsnames]{xcolor} % for coloring text
\definecolor{primaryColor}{RGB}{34, 85, 60} % define primary color
\usepackage{enumitem} % for customizing lists
\usepackage{fontawesome5} % for using icons
\usepackage{amsmath} % for math
\usepackage[
    pdftitle={Connor Haynes' CV},
    pdfauthor={Connor Haynes},
    colorlinks=true,
    urlcolor=primaryColor
]{hyperref} % for links, metadata and bookmarks
\usepackage[pscoord]{eso-pic} % for floating text on the page
\usepackage{calc} % for calculating lengths
\usepackage{bookmark} % for bookmarks
\usepackage{lastpage} % for getting the total number of pages
\usepackage{changepage} % for one column entries (adjustwidth environment)
\usepackage{paracol} % for two and three column entries
\usepackage{ifthen} % for conditional statements
\usepackage{needspace} % for avoiding page brake right after the section title
\usepackage{iftex} % check if engine is pdflatex, xetex or luatex
 
% Ensure that generate pdf is machine readable/ATS parsable:
\ifPDFTeX
    \input{glyphtounicode}
    \pdfgentounicode=1
    \usepackage[T1]{fontenc}
    \usepackage[utf8]{inputenc}
    \usepackage{lmodern}
\fi
 
\usepackage{charter}
 
% Some settings:
\raggedright
\AtBeginEnvironment{adjustwidth}{\partopsep0pt} % remove space before adjustwidth environment
\pagestyle{empty} % no header or footer
\setcounter{secnumdepth}{0} % no section numbering
\setlength{\parindent}{0pt} % no indentation
\setlength{\topskip}{0pt} % no top skip
\setlength{\columnsep}{0.15cm} % set column seperation
\pagenumbering{gobble} % no page numbering
 
\titleformat{\section}{\color{primaryColor}\needspace{4\baselineskip}\bfseries\large}{}{0pt}{}[\color{primaryColor}\vspace{1pt}\titlerule]
 
\titlespacing{\section}{
    % left space:
    -1pt
}{
    % top space:
    0.3 cm
}{
    % bottom space:
    0.2 cm
} % section title spacing
 
\renewcommand\labelitemi{$\vcenter{\hbox{\small$\bullet$}}$} % custom bullet points
\newenvironment{highlights}{
    \begin{itemize}[
        topsep=0.10 cm,
        parsep=0.10 cm,
        partopsep=0pt,
        itemsep=0pt,
        leftmargin=0 cm + 10pt
    ]
}{
    \end{itemize}
} % new environment for highlights
 
 
\newenvironment{highlightsforbulletentries}{
    \begin{itemize}[
        topsep=0.10 cm,
        parsep=0.10 cm,
        partopsep=0pt,
        itemsep=0pt,
        leftmargin=10pt
    ]
}{
    \end{itemize}
} % new environment for highlights for bullet entries
 
\newenvironment{onecolentry}{
    \begin{adjustwidth}{
        0 cm + 0.00001 cm
    }{
        0 cm + 0.00001 cm
    }
}{
    \end{adjustwidth}
} % new environment for one column entries
 
\newenvironment{twocolentry}[2][]{
    \onecolentry
    \def\secondColumn{#2}
    \setcolumnwidth{\fill, 4.5 cm}
    \begin{paracol}{2}
}{
    \switchcolumn \raggedleft \secondColumn
    \end{paracol}
    \endonecolentry
} % new environment for two column entries
 
\newenvironment{threecolentry}[3][]{
    \onecolentry
    \def\thirdColumn{#3}
    \setcolumnwidth{, \fill, 4.5 cm}
    \begin{paracol}{3}
    {\raggedright #2} \switchcolumn
}{
    \switchcolumn \raggedleft \thirdColumn
    \end{paracol}
    \endonecolentry
} % new environment for three column entries
 
\newenvironment{header}{
    \setlength{\topsep}{0pt}\par\kern\topsep\centering\linespread{1.5}
}{
    \par\kern\topsep
} % new environment for the header
 
\newcommand{\placelastupdatedtext}{% \placetextbox{<horizontal pos>}{<vertical pos>}{<stuff>}
  \AddToShipoutPictureFG*{% Add <stuff> to current page foreground
    \put(
        \LenToUnit{\paperwidth-2 cm-0 cm+0.05cm},
        \LenToUnit{\paperheight-1.0 cm}
    ){\vtop{{\null}\makebox[0pt][c]{
        \small\color{gray}\textit{Last updated in October 2026}\hspace{\widthof{Last updated in October 2026}}
    }}}%
  }%
}%
 
% save the original href command in a new command:
\let\hrefWithoutArrow\href

% new command for external links:


\begin{document}
    \newcommand{\AND}{\unskip
        \cleaders\copy\ANDbox\hskip\wd\ANDbox
        \ignorespaces
    }
    \newsavebox\ANDbox
    \sbox\ANDbox{$|$}

    \begin{onecolentry}
        \fontsize{25 pt}{25 pt}\selectfont \textbf{Connor J. Haynes}

        \vspace{5 pt}

        \normalsize Department of Mathematics, University of Western Ontario, London, ON, Canada.

        \vspace{2 pt}

        \textbf{Email:} \hrefWithoutArrow{mailto:chayne8@uwo.ca}{chayne8@uwo.ca}
        \hspace{10 pt}
        \textbf{Website:} \hrefWithoutArrow{https://connorjhaynes.github.io/}{https://connorjhaynes.github.io/}
    \end{onecolentry}

    \vspace{5 pt - 0.3 cm}

    \section{Mathematical interests}
        \begin{onecolentry}
          Broadly: Combinatorics, geometry, and algebra with an focus on computational approaches and the application of these tools to problems in the sciences.\\
          \vspace{0.1cm}
          Specifically: Matroid theory, polyhedral geometry, combinatorics of generic sparse systems, Combinatorial Hodge Theory, geometric tomography, and the application of these tools to problems in the sciences.
        \end{onecolentry}

    \section{Education}
        \begin{twocolentry}{
            Aug 2026 – Aug 2028
        }
            \textbf{MSc in Mathematics}, University of Western Ontario
        \end{twocolentry}
        \vspace{0.1cm}
        \begin{onecolentry}
            Advisor: Taylor Brysiewicz\\
            Thesis title (in progress): \emph{Algebraic matroids of generic sparse systems}
        \end{onecolentry}
        \vspace{0.2cm}

        \begin{twocolentry}{
            May 2022 – May 2025
        }
            \textbf{BSc in Mathematics}, Georgia Institute of Technology
        \end{twocolentry}
        \vspace{0.1cm}

        \section{Writings}
        \begin{onecolentry}
            \textbf{Algebraic Curves}
        \end{onecolentry}
        \vspace{0.05cm}
        \begin{twocolentry}{
            2026
        }
            Solutions to Exercises from Fulton's \emph{Algebraic Curves} (\href{https://connorjhaynes.github.io/fultonsolns.pdf}{link})
        \end{twocolentry}
        \vspace{0.2cm}
        \begin{onecolentry}
            \textbf{Commutative Ring Theory}
        \end{onecolentry}
        \vspace{0.05cm}
        \begin{twocolentry}{
            2026
        }
            Solutions to Exercises from Matsumura's \emph{Commutative Ring Theory} (\href{https://connorjhaynes.github.io/matsumurasolns.pdf}{link})
        \end{twocolentry}
        \vspace{0.2cm}
        \begin{onecolentry}
            \textbf{Introduction to the Orlik-Solomon algebra}
        \end{onecolentry}
        \vspace{0.05cm}
        \begin{twocolentry}{
            2025
        }
            Expository notes on Orlik-Solomon algebras in geometry and combinatorics (\href{https://connorjhaynes.github.io/orliksolomon.pdf}{link})
        \end{twocolentry}
        \vspace{0.2cm}
        \begin{onecolentry}
            \textbf{Moduli spaces - A collection of interesting examples}
        \end{onecolentry}
        \vspace{0.05cm}
        \begin{twocolentry}{
            2025
        }
            Expository notes on moduli spaces targeted towards undergraduates (\href{https://connorjhaynes.github.io/6441paper.pdf}{link})
        \end{twocolentry}
        
        \section{Workshops}
        \begin{onecolentry}
            \textbf{University of Notre Dame}
        \end{onecolentry}
        \vspace{0.05cm}
        \begin{twocolentry}{
            June - July 2026
        }
            Center for Mathematics at Notre Dame: Thematic Program in Algebraic Combinatorics and Applications
        \end{twocolentry}
        \vspace{0.2cm}
        \begin{onecolentry}
            \textbf{Georgia Institute of Technology}
        \end{onecolentry}
        \vspace{0.05cm}
        \begin{twocolentry}{
            May 2026
        }
            Future of Macaulay2 Workshop
        \end{twocolentry}
        \vspace{0.2cm}
        \begin{onecolentry}
            \textbf{University of Notre Dame}
        \end{onecolentry}
        \vspace{0.05cm}
        \begin{twocolentry}{
            June - July 2025
        }
            Center for Mathematics at Notre Dame: Thematic Program in Discrete Groups in Topology and Algebraic Geometry
        \end{twocolentry}
        \vspace{0.2cm}
        \begin{onecolentry}
            \textbf{Louisiana State University}
        \end{onecolentry}
        \vspace{0.05cm}
        \begin{twocolentry}{
            May 2025
        }
            RTG Summer Workshop: Matroids and Hyperplane Arrangements
        \end{twocolentry}

        \section{Teaching}
        \begin{onecolentry}
          \textbf{Graduate Teaching Assistant}, University of Western Ontario
        \end{onecolentry}
        \vspace{0.05cm}
        \begin{twocolentry}{
	    2026
          }
          MATH 1228 - Methods of Finite Mathematics
        \end{twocolentry}
        \vspace{0.2cm}
        \begin{onecolentry}
          \textbf{Student Assistant}, Georgia Institute of Technology
        \end{onecolentry}
        \vspace{0.05cm}
        \begin{twocolentry}{
	    2025
          }
          MATH 3012 - Applied Combinatorics
        \end{twocolentry}
        \vspace{0.2cm}
        \begin{onecolentry}
          \textbf{Undergraduate Teaching Assistant}, Georgia Institute of Technology
        \end{onecolentry}
        \begin{twocolentry}{
            2025
          }
          MATH 1552 - Integral Calculus\\
        \end{twocolentry}
        \vspace{0.05cm}
        \begin{twocolentry}{
	    2024
          }
          MATH 1554 - Linear Algebra\\
        \end{twocolentry}
        
        \section{Service}
        \begin{onecolentry}
          \textbf{UWO Directed Reading Program} - An Introduction to Hyperbolic Geometry
        \end{onecolentry}
        \vspace{0.05cm}
        \begin{twocolentry}{
	    2026
          }
          Student: Evan Korabline
        \end{twocolentry}
        
        \section{Programming languages}
        \begin{onecolentry}
          \texttt{Julia}, \texttt{Macaulay2}, and \texttt{Sage} with emphasis on problems in algebraic geometry and combinatorics.\\
          \vspace{0.1cm}
          \texttt{Python} with emphasis on modelling and functional scripting.
        \end{onecolentry}
\end{document}
